\section{Scope, risks, and exit criteria}

In scope: constant-product AMM, external CEX price, provisional representative agent
roles, a traceable local discrete map, E0--E3, reproducibility, synthetic experiments,
and limited comparison with market facts.

Out of scope: profitable forecasting, operational exchanges, full limit-order-book
replicas, blockchain consensus, advanced MEV, lending, bridges, reinforcement
learning, and general guarantees across all markets. The present MVP does not establish
empirical validity for the map. In its first held-out check, its Schur label matched the
majority contraction outcome in only 3 of 8 groups. The main risk is that the reduced
map may not approximate the ABM outside a local region. Mitigation: review agent rules,
separate model selection from confirmation, record configurations, compare ablations,
and report disagreement as a result. Other risks include overparameterization, seed
dependence, simplified liquidity, confusion between local and global stability, and
calibration bias.

\section{Discussion and limits to inference}

The analysis must distinguish three levels: mechanical properties derived from the
ideal AMM, local diagnostics from the reduced map, and patterns from a stochastic
discrete ABM that remain subject to validation. The present map classified only 3 of 8
held-out majority outcomes correctly under a narrow change in seed and shock size; its
Schur condition does not guarantee stability in the ABM or in real markets. Nor does
the selection run replace evaluation with public data. Further work must report the
regimes in which the approximation reduces error and those in which it fails because
of saturation, capital limits, persistent shocks, or liquidity changes.

\section{Conclusions and future work}

The project now has a reproducible deterministic mechanism runner, a paired
stochastic selection model, and a held-out reduced-map diagnostic. The deterministic
checks show that arbitrage capacity changes correction speed and lower reserves
increase fixed-sale impact in the configured CPMM. The stochastic comparison shows
that the adaptive-LP effect depends on shock direction. The first H1c holdout agrees in
only 3 of 8 majority-outcome groups, so the current map is not a reliable classifier.
The next phase is to finalize the agent roles, broaden the operating regions, preregister
confirmation criteria, and evaluate against a bounded public dataset. None of the
current results demonstrates usefulness in a real market.
